\section{Theorem M*: the macroscopic motion of the work heads}\label{sec:mstar}

Discrete Isoperimetry (\cref{cor:isoperimetry,cor:isoperimetry-trace}) says that, for every
partition of time into small classes with few boundaries, $\Omega(n)$ revisited blocks must link
different classes. This section passes to the limit $n\to\infty$. Rescaled, the head trajectories of
a linear-time transposer converge (along subsequences) to $1$-Lipschitz paths, and the counting
inequality becomes an integral inequality~$(\star)$ for these paths: in every partition of
macroscopic time into pieces of measure at most $\theta$, the heads must return to the same
positions in different pieces with total multiplicity at least $1-2\theta$. In particular the
heads cannot all sweep monotonically: some head makes a macroscopic U-turn, and together the heads
waste at least one unit of range.

The argument uses the cut inequality only, not the machine; we state it for arbitrary walks
(Theorem~M* for walks) and obtain the machine forms (oblivious, trace classes, adaptive, robust,
non-uniform) and, in \cref{sec:codes}, the form for network codes as corollaries.

\subsection{Rescaled walks and their limits}

Fix $G\ge1$. At size $n$ the block size is $r=\lceil nb(n)/G\rceil$ (\lean{TSDiagram.blockSize}), so a
block holds a row. For a walk $h:\N\to\Z$ the \emph{rescaled walk} is
\[
  P^{(n)}(u)=\frac{1}{n}\,\Bigl\lfloor \frac{h(\lfloor un\rfloor r)}{r}\Bigr\rfloor ,
\]
the block at the first step of slot $\lfloor un\rfloor$, in units of $n$ blocks
(\lean{blockTraj}). For a machine, $P^{(n)}_t$ is the rescaled trajectory of work head $t$ on the
all-zero input (\lean{rescaledTraj}), or on an input $w$ (\lean{rescaledTrajOf}). If $T\le cL$,
macroscopic time runs over $[0,c_0]$ with $c_0=cG+1$.

\begin{lemma}[Compactness and regularity]\label{lem:walk-basics}
\formal{Tape/Statements/Coding.lean}{blockTraj_limit_lipschitz, exists_subseq_blockTraj, numSlots_le_of, frozen_V2, pairRes_congr}
Let $n_i\to\infty$ and let $h_i$ be walks with unit steps and block sizes $r_i\ge1$.
\begin{enumerate}[label=(\alph*)]
\item Every uniform limit on $[0,c_0]$ of the rescaled walks is $1$-Lipschitz.
\item If the walks start at $0$, then for any $k$ walks per size some subsequence of the rescaled
walks converges uniformly on $[0,c_0]$.
\item If $T\le cn^2b$ and a block holds a row ($nb\le rG$), there are at most $(cG+1)n$ slots.
\item[(V2$'$)] For a walk frozen at time $T$, every block strictly between its blocks at the first
steps of two slots $y_1\le y_2$ is pair-resident at some slot of $[y_1,y_2]$; and the pair residency
depends only on the walk before $T$.
\end{enumerate}
\end{lemma}

\begin{outline}
(a) A unit-step walk moves at most $r$ cells per slot, that is at most one block per $1/n$ of
macroscopic time, up to rounding that vanishes in the limit. (b) Arzelà--Ascoli in the form of a
diagonal argument at rational times (a countable product of compact intervals is compact and
metrizable) together with asymptotic equicontinuity; the classical statement does not apply
directly because the rescaled walks are step functions. (c) Arithmetic. (V2$'$) is
\cref{lem:residency}(V2) for frozen walks.
\end{outline}

\begin{lemma}[Machine instances]\label{lem:machine-walks}
\formal{Tape/Statements/Core.lean}{pairRes_V2, numSlots_le_horizon, rescaledTraj_limit_lipschitz, exists_subseq_tendstoUniformlyOn}
For every machine $M$, time constant $c$ and $n_i\to\infty$: some subsequence of the rescaled
work-head trajectories converges uniformly on $[0,c_0]$, every such limit is $1$-Lipschitz, and if
$T\le cL$ there are at most $c_0n$ slots. If $M$ halts at $T$ on $w$, its work heads satisfy
(V2$'$).
\end{lemma}

\begin{outline}
Instances of \cref{lem:walk-basics}.
\end{outline}

\subsection{Levels, labellings and the inequality \texorpdfstring{$(\star)$}{(*)}}

\begin{definition}[Levels and crossings]\label{def:levels}
\formal{Tape/DefsM.lean}{levelSet, ConsecIn, crossPairs, levelCross, levelExcess, FinUnionIcc, PiecewiseConstOn}
Let $P:\R\to\R$ and $c_0>0$. The \emph{level set} of $p$ is $\Lambda(p)=P^{-1}(p)\cap[0,c_0]$. An
\emph{excised set} $S$ is a finite union of closed intervals, and a \emph{labelling}
$\beta:\R\to\N$ is piecewise constant off $S$: constant on every interval of $[0,c_0]\setminus S$
that avoids a fixed finite set. For a level $p$ with $\Lambda(p)$ finite, $\kappa(p)$ is the number
of consecutive pairs $s<s'$ of $\Lambda(p)$ with $s,s'\notin S$ and $\beta(s)\ne\beta(s')$; for
$\Lambda(p)$ infinite, $\kappa(p)=0$. The \emph{excess} of the level is
$\#\Lambda(p)-1$ (truncated at $0$, and $\infty$ for an infinite level set).
\end{definition}

So $\kappa(p)$ counts the returns of the path to height $p$ that cross from one label class to
another. For $k$ paths $P_t$ on $[0,c_0]$, a number $\theta\ge0$, an excised set $S$ and a
labelling $\beta$ whose classes off $S$ have measure at most $\theta$, the inequality is
\[
  (\star)\qquad
  1\;\le\;\sum_{t<k}\int\kappa_t(p)\,dp\;+\;2\theta\;+\;2k\,\bigl|S\cap[0,c_0]\bigr| .
\]

\begin{definition}[The cut inequality]\label{def:cutineq}
\formal{Tape/DefsM.lean}{CutIneq}
Walks $h_t$ ($t<k$) of length $T$ satisfy the \emph{cut inequality} at size $n$ with block size
$r$, $G$ bits per cell, slack $D$ and boundary cost $s$, if for every labelling of the
$N=\lceil T/r\rceil$ slots with at most $m$ slots per label and $\Mb$ label changes,
\[
  (n\dotminus\Mb)(n\dotminus2m)\,b(n)\;\le\;(\Mb+1)D+\Mb\,s+rG\,\bigl(3k\Mb+\chicross\bigr),
\]
where $\chicross$ counts the chords of the pair residencies whose endpoints carry different labels, and
$x\dotminus y=\max(x-y,0)$.
\end{definition}

For the work heads of a machine on a trace class it follows from \cref{cor:isoperimetry-trace} (since
$\clog{(g+1)}\le G$), with $D$ the co-entropy of the class and $s=\clog{(q+1)}+(k+1)\clog{(2T+1)}$; the formal
proofs obtain it through the class code and the cut-set bound for codes (\cref{lem:class-code},
\cref{cor:isoperimetry-code}).

\subsection{Theorem M* for walks}

\begin{theorem}[Theorem M* for walks]\label{thm:mstar-walk}
\formal{Tape/Statements/Coding.lean}{theoremMStar_walk, theoremMStar_level_walk}
Let $n_i\to\infty$ and, eventually, let $k$ walks $h_{i,t}$ have unit steps before $T_i\le cL_i$
and satisfy the cut inequality at size $n_i$ with block size $\lceil n_ib(n_i)/G\rceil$, slack
$D_i$ and boundary cost $s_i$, where $(D_i+s_i)/L_i\to0$. Let $c_0=cG+1$ and consider the rescaled
walks frozen at $T_i$.
\begin{enumerate}[label=(\alph*)]
\item If the walks start at $0$, some subsequence converges uniformly on $[0,c_0]$.
\item Every subsequential uniform limit $(P_t)_{t<k}$ consists of $1$-Lipschitz paths and
satisfies $(\star)$ for every $\theta$, $S$ and $\beta$. Moreover
\[
  \sum_t\int(\#\Lambda_t(p)-1)\,dp\;\ge1,\qquad
  1+\sum_t\bigl|P_t([0,c_0])\bigr|\;\le\;k\,c_0,
\]
and the paths are not all monotone.
\end{enumerate}
\end{theorem}

\begin{outline}
The discrete part works at one size. From $S$ and $\beta$ build a slot labelling: the label of a
slot off $S$ is determined by $\beta$, and slots in $S$ are cut into short pieces. Its label changes
and class sizes are bounded by counting good and bad slots. The \emph{chord split}
(\lean{crossChords_split}) bounds the crossing chords by the label changes times $(\lceil\eta n\rceil+1)$, plus
twice the slots in $S$, plus the \emph{long} crossing chords, those spanning at least $\eta n$
slots. Long chords are counted by levels (\lean{levelChords_limsup_le}): for a good level $p$, the
long crossing chords at block $\lfloor pn\rfloor$ inject, for large $n$, into the pairs counted by
$\kappa(p)$, by uniform closeness, (V1) and (V2$'$). Reverse Fatou then gives
$\limsup_n\#\{\text{long crossing chords}\}/n\le\int\kappa(p)\,dp$, with an explicit integrable
majorant and no measurability assumption on $\kappa$. If $(\star)$ failed, the cut inequality would
fail at a single large size for explicit $\eta$ and $\delta$; the arithmetic is isolated in one
lemma. The remaining claims are the labelling-free consequences of \cref{thm:excess} below.
\end{outline}

The interval labellings and excised sets avoid the area formula and Borel measurability of first
returns: only finitely many breakpoints occur.

\subsection{Labelling-free consequences}

\begin{theorem}[Excess and range]\label{thm:excess}
\formal{Tape/Statements/Core.lean}{levelCross_le_levelExcess, one_le_excess_of_star, excess_add_range_le, range_add_one_le_of_star, lintegral_levelExcess_eq_zero_of_monotone, not_all_monotone_of_star}
For every path, excised set and labelling, $\kappa(p)\le\#\Lambda(p)-1$. Let $P_t$ ($t<k$) satisfy
$(\star)$ for every $\theta\in(0,\tfrac12)$, $S$ and $\beta$. Then:
\begin{enumerate}[label=(\alph*)]
\item the total excess multiplicity is at least one:
$\sum_t\int(\#\Lambda_t(p)-1)\,dp\ge1$;
\item if the $P_t$ are $1$-Lipschitz on $[0,c_0]$, then
$\int(\#\Lambda(p)-1)\,dp+\abs{P([0,c_0])}\le c_0$ for each $t$, and hence
$1+\sum_t\abs{P_t([0,c_0])}\le kc_0$;
\item a monotone path has total excess $\int(\#\Lambda(p)-1)\,dp=0$, so the $P_t$ are not all monotone.
\end{enumerate}
\end{theorem}

\begin{outline}
For~(a), take $S=\emptyset$ and the labelling $\max(0,\lfloor mu\rfloor)$ with $\theta=1/m$, and let
$m\to\infty$. Part~(b) is the easy half of Banach's indicatrix~\cite{Banach1925} for $1$-Lipschitz
paths, proved by dyadic counting with the range split off. For~(c), a level with two preimages
under a monotone path is a value on an interval of constancy, and there are countably many such
values.
\end{outline}

\begin{lemma}[Analysis tools]\label{lem:analysis}
\formal{Tape/Statements/Core.lean}{countable_localExtr_values, levelSet_encard_majorant, lintegral_encard_levelSet_le, ae_finite_levelSet, levelChords_limsup_le, crossChords_split}
For every $P:\R\to\R$ the values at local extrema form a countable set. For $P$ that is
$K$-Lipschitz on $[0,c_0]$, $\#\Lambda(p)$ has a measurable majorant with integral at most $Kc_0$;
hence $\int\#\Lambda(p)\,dp\le Kc_0$ and almost every level set is finite. Moreover:
\begin{enumerate}[label=(\alph*)]
\item (Chord split.) Let a residency $R$ satisfy (P2), let $\eta\in\R$, and let a slot labelling
$\lambda$ satisfy $\beta(y/n)\ne\beta(y'/n)$ whenever $y/n,y'/n\notin S$ and $\lambda(y)\ne\lambda(y')$.
Then the chords of $R$ with target $<N$ whose ends carry different labels number at most
$\Mb(\lceil\eta n\rceil+1)+2\,\#\{y<N:y/n\in S\}+\#\mathcal X$, where $\mathcal X$ is the set of
long chords $x\to z$: $z<N$, $z-x\ge\eta n$, $x/n,z/n\notin S$, $\beta(x/n)\ne\beta(z/n)$.
\item (Long chords by levels.) Let $\eta>0$, let $P$ be $1$-Lipschitz on $[0,c_0]$, and for large $i$ let
$W_i:\N\to\Z$, residencies $R_i$ and $N_i\le c_0n_i$ satisfy $\abs{W_i(y)-\sigma}\le D$ for
$\sigma\in R_i(y)$ and (V2$'$), with $u\mapsto W_i(\lfloor un_i\rfloor)/n_i\to P$ uniformly on $[0,c_0]$.
Then $\limsup_i\#\mathcal X_i/n_i\le\int\kappa(p)\,dp$ for every excised set $S$ and labelling $\beta$.
\end{enumerate}
\end{lemma}

\begin{outline}
Each value at a local extremum is the minimum or maximum of $P$ over a rational interval. The
majorant is $\liminf_mN_m(p)$, where $N_m(p)$ counts the dyadic subintervals $I$ of order $m$ such that $p$
lies within $K\abs{I}/2$ of the value of $P$ at the centre of $I$ (an interval of values containing $P(I)$).
Each $N_m$ is measurable with integral $Kc_0$; once $c_02^{-m}$ is below the gaps of a finite subset of
$\Lambda(p)$, $N_m(p)$ is at least its size, so $\liminf_mN_m\ge\#\Lambda$ at every level, and Fatou's lemma
bounds the integral. For~(a), a crossing chord shorter than $\eta n$ has its source among the
$\lceil\eta n\rceil$ slots before a label change; a chord touching $S$ is charged to its source or target slot
in $S$; by (P2) both charges are injective, and the remaining crossing chords are long. Part~(b) is the level
count in the outline of \cref{thm:mstar-walk}.
\end{outline}

\subsection{Machine forms}

\begin{theorem}[Theorem M*, oblivious machines]\label{thm:mstar}
\formal{Tape/Statements/Core.lean}{theoremMStar_level, theoremMStar_level_sub, theoremMStar_level_subseq, theoremMStar_excess, no_stationary_limits}
Let $n_i\to\infty$ and let $M$ be, eventually, oblivious on $\mathrm{Mat}(n_i,b(n_i))$ with halting
time $T_i\le cL_i$ and transpose. If the rescaled work-head trajectories converge uniformly on
$[0,c_0]$ to $P_t$, then the $P_t$ satisfy $(\star)$ for every $\theta$, $S$, $\beta$ (equivalently,
$\sum_t\int\kappa_t\ge1-2\theta-2k\abs{S\cap[0,c_0]}$ for $\theta\ge0$), and the conclusions of
\cref{thm:excess}. A convergent subsequence always exists. In particular the limits cannot all be
identically $0$.
\end{theorem}

\begin{theorem}[Theorem M* for trace classes and adaptive machines]\label{thm:mstar-adaptive}
\formal{Tape/Statements/Adaptive.lean}{rescaledTrajOf_close, tendstoUniformlyOn_of_mem, mstar_discrete_class_on, theoremMStar_level_class, theoremMStar_level_class_on, theoremMStar_class_nonuniform, theoremMStar_level_adaptive, no_stationary_limits_class_on}
Let $n_i\to\infty$.
\begin{enumerate}[label=(\alph*)]
\item Members of one trace class have rescaled trajectories within $4/n$ of each other, so limits
along trace classes do not depend on the representatives.
\item Along trace classes $F_i$ with halting times $T_i\le cL_i$ and co-entropies $D_i=o(L_i)$ on which $M$
is correct, the rescaled trajectories of the representatives satisfy the conclusion of \cref{thm:mstar-walk}
(the labelling-free parts by \cref{thm:excess}); the underlying finite-$n$ inequality holds at every size. The same holds for non-uniform machines
$M_i$ with $k$ tapes, $G$ bits per cell and $o(L_i)$ state bits.
\item If every matrix input halts within $cL_i$ steps and $M$ transposes, then the popular classes
(\cref{thm:popular}) have $D_i\le Cn_ib(n_i)$, with $C=C(M,c)$ the constant of \cref{thm:popular}
(\lean{traceConst}), a convergent subsequence of the representatives'
trajectories exists, and every subsequential limit is $1$-Lipschitz and satisfies $(\star)$.
\end{enumerate}
\end{theorem}

\begin{theorem}[Robust and non-uniform M*]\label{thm:mstar-robust}
\formal{Tape/Statements/Adaptive.lean}{theoremMStar_level_robust, theoremMStar_level_nonuniform, theoremMStar_excess_robust, theoremMStar_excess_nonuniform}
Let $n_i\to\infty$ and suppose that, eventually, $M$ halts within $cL_i$ steps with the
transposed output on every input of a set $W_i$ of co-entropy $m_i$, with $m_i/L_i\to0$; inputs
outside $W_i$ are unconstrained. Then the popular trace classes inside $W_i$ have co-entropy at most
$m_i+Cn_ib(n_i)$, convergent subsequences of the representatives' trajectories exist, and every
subsequential limit is $1$-Lipschitz, satisfies $(\star)$, has total excess at least $1$ and total
range at most $kc_0-1$, and its paths are not all monotone; here $C=C(M,c)$ is the constant of
\cref{thm:popular}. The same holds for machines $M_i$ depending on $i$, with $k$ work tapes, $G$ bits per cell
and $o(L_i)$ state bits (then $C$ depends only on $k$, $G$ and $c$).
\end{theorem}

\begin{outline}[Outline of the formal proofs of \cref{thm:mstar,thm:mstar-adaptive,thm:mstar-robust}]
All machine forms reduce to Theorem~M* for code-carrying diagrams (\cref{thm:mstar-diagram}),
hence to Theorem~M* for walks. For a trace class, the class code on the representative's diagram
(\cref{lem:class-code}) carries the state and head positions as
$\clog{(q+1)}+(k+1)\clog{(2T+1)}=o(L)$ extra bits per memory edge, and the cut-set bound for codes
gives the cut inequality. The oblivious machine is the case of one class with $D=0$
(\cref{lem:trace-obl}); the adaptive form is the robust form with $W_i$ all inputs; the robust form
extracts the popular class inside $W_i$ (\cref{thm:popular}, in its form within a set of inputs,
\lean{exists_popular_traceClass_on}). Representative independence holds
because the work-head blocks of two members of a class differ by at most $2$ before halting and $4$
after.
\end{outline}

\begin{corollary}[Theorem M* under (T)]\label{cor:mstar-T}
\formal{Tape/Statements/Adaptive.lean}{adaptive_hyps_of_linearTime, linearTimeTransposable_popular_limits}
If (T) holds, then some machine $M$ and constant $c$ have the following property along every
$n_i\to\infty$: there are trace classes of co-entropy $o(L_i)$ (the popular traces), a convergent
subsequence of their representatives' rescaled trajectories, and every subsequential limit is
$1$-Lipschitz and satisfies $(\star)$ on $[0,(2c)G+1]$.
\end{corollary}

\begin{outline}
A machine meeting \cref{def:T} with constant $c$ halts within $2cL$ steps on every matrix input of every
size $n\ge1$ and transposes (\cref{lem:bridge}); apply \cref{thm:mstar-adaptive}(c).
\end{outline}

So a linear-time transposer must move its work heads in a macroscopically folding way. The
condition is far from sufficient: it constrains only the macroscopic shape of the head motion. It does
exclude simple zig-zags: if every head follows the same triangle wave $P\ge0$ of slope $\pm1$, the labelling
$\beta(u)=\lfloor mP(u)\rfloor$ with $m$ large has classes of measure at most $\tfrac14$ and keeps every level
set inside one class, so $\kappa_t\equiv0$ and $(\star)$ fails at $\theta=\tfrac14$, $S=\emptyset$. Its value
is as a test that any candidate construction must pass, and as the limit object behind the coding results of
\cref{sec:codes}.

\begin{proposition}[$(\star)$ is a real constraint]\label{prop:star-fails}
\formal{Tape/Statements/Core.lean}{theoremMStar_conclusion_can_fail}
For the constant path $P=0$, $S=\emptyset$, $\beta(u)=\max(0,\lfloor4u\rfloor)$ and $\theta=\tfrac14$, the
hypotheses on $S$, $\beta$ and $\theta$ hold and $(\star)$ fails, for every $k$ and $c_0$.
\end{proposition}

\begin{outline}
$\kappa\le\#\Lambda-1$, and the excess of a constant path vanishes almost everywhere.
\end{outline}
